\chapter{Espaces mesurables et probabilité}
\section{Tribu}

\begin{definition}[Tribu] 
    Soit $\mathcal{A} \in \mathcal{P}(E)$. On dit que $\mathcal{A}$ 
    est une tribu si et seulement si :

    \begin{enumerate}
        \item $E \in \mathcal{A}$
        \item $\forall A \in \Acal \implies A^{c} \in \Acal $
        \item $\forall (A_n) \in \Acal^\mathbf{N}$ ; $\bigcup_{n\in \mathbf{N}} A_n \in \Acal$
    \end{enumerate}

\end{definition}

\begin{remarque}
    Par la relation d'inclusion sur $\Parts{E}$ alors ${\emptyset,E}$ 
    est  la plus petite tribu sur E, dite grossière. Et $\Parts{E}$ la plus grande
\end{remarque}

\begin{propriété} 
    \begin{enumerate}
        \item  On peut remplacer (i) par $ \emptyset \in \Acal$
        \item $ \emptyset \in  \Acal \implies \Acal$ est stable par intersectiopn et union finie
        \item Soit $A,B \in \Acal$ alors $ B\setminus A \in \Acal$ et $ B \Delta A \in \Acal$
    \end{enumerate}
\end{propriété}

\begin{remarque}
    Une union de tribu n'est pas en général une tribu.
\end{remarque}

\begin{exemple}
$ E = \{1;2;3\}$ \\
$\Acal_1 = \{ \emptyset ; \{1;2;3\}; \{1\} ; \{2;3\} \}$ \\
$\Acal_2 = \{ \emptyset ; \{1;2;3\}; \{3\} ; \{1;2\} \}$ \\
$ \{1;2\} , \{2;3\} \in    \Acal_1 \cup \Acal_2 $ \\
Mais $ \{1;2\} \cap \{2;3\} = \{2\} \Acal_1 \cup \Acal_2$
\end{exemple}

\begin{propriété}
    Une intersection quelconque de tribu sur E est une tribu
\end{propriété}

\begin{definition}[Tribu engendré]
\end{definition}

\begin{exemple}
\end{exemple}

\begin{remarque}
\end{remarque}

\begin{definition}[Tribu produit]
\end{definition}

\begin{definition}[Tribu image reciproque]
\end{definition}

\begin{definition}[Tribu trace]
\end{definition}

\begin{propriété}
\end{propriété}

\begin{remarque}
\end{remarque}